\section*{Kapitel 10 --- Drivrutinen}
\addcontentsline{toc}{section}{Kapitel 10 --- Drivrutinen}

\solution{ex:10:1}{Bryt koden med flit}
Uppgiften ber om kodändringar; utfallen beskrivs här.
\ans{a} Triggern skrivs innan dataregistren, så kontrollern läser föregående frames innehåll.
På bussen syns en nod som ligger exakt en frame efter --- den skickar alltid det du bad om
\emph{förra} gången.
\ans{b} Databyte 4--7 hamnar i \code{TX_DATA_LO}. En frame med \code{dlc = 2} och data
\code{AA BB} skickas som två nollbytes; mottagaren får rätt längd och fel innehåll.
\ans{c} \code{AA BB CC DD} kommer fram som \code{DD CC BB AA}. Framen kommer fram, med rätt
längd --- bara bakvänd.
\ans{d} Felbiten går aldrig ned, så \code{hasError()} returnerar \code{true} för alltid efter
första felet och all felhantering ovanför slutar fungera.
\ans{e} \code{STATUS} bit 1 står kvar. Varje \code{receive()} returnerar \code{true} med samma
frame, i en oändlig loop, och ingen ny frame kan tas emot.
\ans{f} Ingenting syns i utskriften, och det är poängen. \code{PrintTransport} visar vilka bytes
som gick ut, och de är korrekta --- felet ligger i vad drivrutinen \emph{lämnar ifrån sig} till
anroparen, alltså i \code{frame.data} ovanför \code{dlc}. Det fångas av ett testfall som levererar
en frame med \code{dlc = 2} efter en med \code{dlc = 8} och kontrollerar att de sex översta
bytesen är nollställda, vilket kräver en dubblare som modellerar hårdvarans beteende --- alltså
\code{BankTransport} i \kapref{ch:testning}, inte den skriptade.
\ans{g} Skrivningar sker mitt i en pågående sändning och skriver över data kontrollern använder.
Utfallet beror på timingen, vilket gör felet sporadiskt.

\solution{ex:10:2}{Räkna transaktioner}
\ans{a} Sex: en \code{STATUS}-läsning och fem skrivningar.
\ans{b} En --- \code{STATUS}-läsningen. Därefter returnerar den \code{false}.
\ans{c} \textbf{Noll.} Valideringen sker före allt annat, så en ogiltig frame kostar inte en enda
transaktion.
\ans{d} En.
\ans{e} Sex: \code{STATUS}, \code{RX_ID}, \code{RX_DLC}, två dataregister och \code{RX_ACK}.

Fråga c är poängen med ordningen: en anropare som råkar skicka ogiltiga frames i en loop belastar
inte bussen alls.

\solution{ex:10:3}{Den oändliga loopen}
Uppgiften ber om kod; utfallen beskrivs här.
\ans{a} Loopen snurrar tills hårdvaran blir ledig, och avslutas sedan normalt. Det är det avsedda
beteendet.
\ans{b} Loopen snurrar för alltid. \code{isValid()} returnerar \code{false} oavsett hur länge man
väntar, och ingenting i systemet kommer att ändra på det.
\ans{c} Validera före loopen och returnera direkt vid en ogiltig frame, och räkna försök i loopen
så att även en hårdvara som aldrig blir ledig ger upp. Se \secref{sec:waiting}.

\solution{ex:10:4}{Designfrågor}
\ans{a} Till exempel en uppräkning som returvärde i stället för \code{bool}: \code{Accepted},
\code{Busy}, \code{Invalid}. Kostnaden för \code{Stub} är liten --- den skulle returnera
\code{Accepted} eller \code{Invalid} --- men kostnaden för \emph{interfacet} är att det växer, och
varje anropare måste hantera fler fall. Avvägningen är att en \code{bool} är enklare att
implementera och svårare att använda rätt.
\ans{b} Ett testfall som levererar en frame med \code{dlc = 8} och därefter en med \code{dlc = 2}.
Utan maskering bär den andra framen kvar sex bytes från den första, och applikationen kan inte
skilja dem från riktig data. På riktig hårdvara krävs dessutom bara att hårdvarugruppen ändrar
när dataackumulatorn nollställs, eller att du lånar ett kort byggt av en annan grupp.
\ans{c} Drivrutinen observerar att \code{STATUS} bit 0 går låg vid triggern och kommer tillbaka
strax därefter --- exakt som vid en lyckad sändning --- och att felflaggan är satt. Den kan
\emph{inte} observera att det var en arbitrering som förlorades snarare än ett stoppfel, en
trasig CRC eller en mottagen fjärrframe, och den kan inte se vilken nod som vann eller vilken
frame som gick igenom i stället.
\ans{d} Att polla oftare, alltså att göra mindre arbete mellan \code{receive()}-anropen. Gränsen
du inte kommer förbi är att en \code{receive()} tar omkring 270~µs medan en frame tar 130~µs: en
sändare som kör kontinuerligt hinner alltid skicka fler frames än du hinner hämta.
